\section{Introduction} 
\label{sec:intro}

In many real-world scenarios, decision-makers are required to make a series of choices over time, with the goal of optimizing a long-range objective without knowing the outcomes in advance. Such sequential decision-making processes (also known as online learning) are prevalent in various fields such as finance~\citep{bell1988game, ban2018machine}, healthcare~\citep{bastani2021efficient, killian2023robust}, and online advertising~\citep{wagner2019digital, choi2020online}. For instance, a financial portfolio manager must decide which assets to invest in daily, balancing potential returns against risks. Similarly, a doctor might choose a treatment plan for a patient, adjusting it based on the response over time. In online advertising, companies continuously select which ads to display to users to maximize engagement and revenue.

There are two dominant models on the nature of the environment generating the data: stochastic or adversarial. In a stochastic environment, the data is assumed to be generated from an unknown, but fixed probability distribution. This assumption allows for the use of statistical methods to predict future outcomes based on historical data, making it possible to optimize decisions under uncertainty. For example, in finance, asset prices might be modeled using stochastic processes to forecast future market trends. Conversely, in an adversarial environment, the data is assumed to be generated by an adversary who aims to thwart the decision-maker’s efforts. This assumption is crucial in scenarios where the decision-maker faces strategic manipulation or worst-case scenarios. Such adversarial models ensure that decision-making strategies remain robust and perform well even when faced with the most challenging and unpredictable conditions, making them ideal for use in high-stakes applications where robust guarantees on performance are important.

Usually, the environment is a-priori assumed to follow one of these two models, and consequently decision-making strategies are designed for optimal performance in the assumed environment. Strategies designed for stochastic environments are typically simple to deploy and can achieve excellent performance for most outcome sequences, they are extremely vulnerable to strategic attacks, where an adversary can exploit their predictability and cause significant losses. On the other hand, worst-case optimal strategies offer stronger guarantees against adversarial manipulation, ensuring robust performance even in the most challenging scenarios. This duality in the strengths and weaknesses of the two aforementioned approaches highlights the need for algorithms that can seamlessly adapt to both benign and adversarial conditions. 
\bluet{
Such a need is also practically relevant; as demonstrated in Bastani et al. [2021b], simple greedy algorithms show remarkable effectiveness on real healthcare datasets. In their contextual bandit experiments across multiple healthcare intervention scenarios, they found that greedy policies (which follow the FTL principle of choosing actions based on empirical performance) achieved near-optimal performance on actual patient data, substantially outperforming worst-case guarantees. Thus, the gap between simple greedy approaches and sophisticated robust algorithms was minimal on realistic datasets, while the greedy approaches maintained significant computational and implementation advantages.
}
However, in healthcare, robustness is crucial to ensure reliable performance across diverse and unpredictable patient responses. This underscores the necessity for “best of both worlds” methods that can leverage the simplicity and efficiency of stochastic strategies while retaining the robustness of worst-case optimal strategies---the problem we study in this paper. \textcolor{blue}{Another example highlighting the need for such methods can be seen in the domain of online advertising. In pay-per-click advertising, platforms must select which ads to display for each user, learning from historical click patterns to optimize revenue. Under normal conditions, simple algorithms that follow historical success patterns work well since user behavior tends to be relatively consistent. However, click fraud introduces adversarial dynamics: malicious actors deploy automated systems to systematically manipulate click patterns. For instance, competitors might use bots to repeatedly search for and ignore a rival's advertisements, artificially making those ads appear ineffective to the platform's learning algorithm. In such environments, algorithms must perform well under typical user interactions while remaining robust against coordinated manipulation campaigns. This exemplifies the practical need for approaches that can leverage the efficiency of simple, data-driven methods during normal operations while maintaining defensive capabilities against strategic attacks.}

There has been a recent surge of interest in developing “best of both worlds” methods for online learning. The typical results give guarantees on the expected regret in a stochastic environment, while maintaining optimal worst case bounds for all instances. The limitation of such a guarantee however, is that generative model assumptions are very difficult to verify. Given a data sequence in hindsight, it may not be clear whether it was drawn from a stochastic environment or not. Furthermore, there may be data sequences that are not drawn from a stochastic environment for which the natural opportunistic algorithm may actually still perform well, yet the typical “best of both worlds” results can only guarantee the worst case bounds. As such we propose a shift in perspective that focuses on guarantees that hold for every possible data sequence, irregardless of how it is generated. In particular, perhaps the truly relevant comparison is the performance of the two benchmark algorithms on the given instance in hindsight. Our goal is to achieve a ``best of both algorithms'' result that guarantees a constant factor of the performance of the best algorithm in hindsight on every data sequence, which we will formally define as ``instance optimality''. This is a stronger guarantee that implies the typical ``best of both worlds'' result.

The key to our proposed solution -- titled \emph{Switching via Monotone Adapted Regret Traces} ($\bob$) -- follows the simple principle of measured optimism: assume the best, prepare for the worst. $\bob$ defaults to playing the opportunistic algorithm $\ftl$ at the beginning, and monitors the performance of the algorithm adaptively. If the opportunistic algorithm seems to perform poorly, then $\bob$ will switch at most once within the time horizon to a given minimax optimal algorithm $\wcalg$, where the switching time is a function of the worst case guarantee on $\wcalg$. Not only is this proposed solution extremely simply to implement, we show that its performance is nearly optimal, guaranteeing a multiplicative factor $e/(e-1) \approx 1.58$ of the minimum regret of $\ftl$ and $\wcalg$, as compared to a lower bound stating that no algorithm can guarantee better than a $1.43$-factor of the minimum regret of $\ftl$ and $\wcalg$. The proof is surprisingly simple, following from a reduction to the online ski rental problem.

\section{Formal problem setup and definition of instance optimality} 
Our work aims to develop algorithms for online learning that are \emph{instance optimal}~\citep{fagin2001optimal},\cite[Chapter $3$]{roughgarden2021beyond} with respect to the stochastic and minimax optimal algorithms for a given setting. This is best motivated via
%a strong but somewhat subtle goal, which we motivate first by considering 
a concrete example:
\begin{example}[Binary Prediction]
We are given bit stream $y^n \defeq y_1,y_2,\ldots,y_n\in\{0,1\}^n$.
At the start of day $t$, before seeing $y_t$, we choose prediction $\widehat{Y}_t \sim \Ber(a_t)$ (for $a_t \in [0,1]$) for the upcoming bit $y_t$, given the history $y^{t-1}$. Our resulting loss on day $t$ is $\ell_t(a_t) = \Prob(\widehat{Y}_t \neq y_t) = |a_t - y_t|$, and our total loss is $L_n(\alg,y^n)\defeq \sum_{t=1}^n\ell_t(a_t)$. The objective is to achieve low \emph{regret} (i.e., additive loss) compared to the loss $L_n(a,y^n)=\sum_{t=1}^n\ell_t(a)$ of the best fixed action $a^* \in[0,1]$ in hindsight. As $a^*$ is the majority in $y^n$ between 0 and 1, it follows that $L_n(a^*, y^n) = \min\left\{\sum_{t=1}^n y_t, n - \sum_{t=1}^n y_t\right\}$. For sequence $y^n \in \{0,1\}^n$, policy $\alg$ incurs regret
\begin{align}
    \Reg(\alg, y^n) \defeq L_n(\alg,y^n) - L_n(a^*,y^n)
     = \textstyle\sum_{t=1}^n |a_t - y_t| - \min_{a\in[0,1]}\textstyle\sum_{t=1}^T|a-y_t|. \label{eq:regdefn}
\end{align}
% In this case, FTL chooses 
% \begin{align} \label{eq:BinFTL}
% \aftl_t = a^*_{t-1} =
% \begin{cases}
%     0 & \text{if Majority}\{y_1,\dotsc,y_{t-1}\} = 0 \\
%     1 & \text{if Majority}\{y_1,\dotsc,y_{t-1}\} = 1 \\
%     \frac{1}{2} & \text{otherwise}.
% \end{cases}
% \end{align}
\end{example}
%\cycomment{The notation above switches between x and y}
Binary prediction goes back to the seminal works of~\cite{blackwell1956analog} and~\cite{hannan1957approximation}. The definition of regret is motivated by the case where $y_t$ is \emph{randomly} generated as i.i.d. Bernoulli$(p)$. If $p$ is known, then the optimal policy is the `Bayes predictor' $a^{\textsf{Bayes}}=\lfloor 2p\rfloor$ (i.e., nearest integer to $p$),
which coincides with hindsight optimal $a^*$ w.h.p. when $p$ is away from $1/2$. When $p$ is unknown, the \emph{stochastic optimal} policy is the \emph{Follow The Leader} or $\ftl$ policy, which sets $a_t=\Majority(y^{t-1})$, i.e. the majority bit amongst the first $t-1$ bits. If there is an equal number of 0s and 1s in $y^{t-1}$, we define $\Majority(y^{t-1})=1/2$, breaking ties uniformly. 

A starting point for online learning is the observation that it is easy to construct a sequence $y^n$ such that $\ftl$ has poor regret: For example, if $y^n = (1,0,1,0,1,0,\ldots)$, i.e., alternate $1$s and $0$s, then the regret of $\ftl$ grows linearly with $n$. In contrast,~\emph{worst-case optimal} online learning policies such as those of Blackwell and Hannan, or more modern versions like Hedge, Multiplicative Weights or FTPL (see~\cite{Cesa-Bianchi--Lugosi2006,slivkins2019introduction}) guarantee regret of $\Theta(\sqrt{n})$ over all sequences. Indeed, for bit prediction, the \emph{exact} minimax optimal policy was established by~\cite{Cover1966}, and this policy (which we refer to as $\cover$) achieves  $\Reg(\cover, y^n) = f_n = \sqrt{\frac{n}{2\pi}}(1+o(1))$ under \emph{any} $y^n \in\{0,1\}^n$. Here $f_n$, the so-called Rademacher complexity of the setting, is a \emph{fixed} function of $n$ that does not depend on sequence $y^n$. For binary prediction, $f_n = \frac{\E|\sum_{t=1}^n Z_t|}{2} \approx \sqrt{\frac{n}{2\pi}}$ where $Z^n \overset{iid}{\sim} \mathrm{Unif}\{1,-1\}$.

While the above discussion seems a convincing endorsement of worst-case online learning algorithms, the situation is more complicated. One problem is that while $\ftl$ has bad regret on certain pathological sequences, on more `realistic' sequences $\ftl$ performs orders of magnitude better than the minimax regret. As an example, with i.i.d. Bernoulli$(p)$ input, $\Reg(\ftl)$ is actually \emph{independent of $n$} (i.e., $O(1)$) as long as $p$ is away from $1/2$. 
%We demonstrate this in~\cref{fig:SMART}$(a)$, where we see $\Reg(\ftl)$ is much lower than $\Reg(\cover)\approx 0.39\sqrt{n}$ unless $p$ is very close to $1/2$. 
This phenomena is known in more general settings~\citep{Huang--Lattimore--Gyorgy--Szapesvari2016}, \bluet{suggesting that for many naturally occurring sequences, FTL may achieve much better performance than worst-case bounds would indicate---see also the discussion in Section~\ref{sec:intro} on positive empirical performance of greedy, $\ftl$-like policies on real data. More broadly, \citep{kotlowski2018minimaxity} shows that $\ftl$ achieves optimal pseudoregret under minimal distributional assumptions in stochastic settings.}
 On the other hand, 
%as~\cref{fig:SMART}$(b,c)$ indicates, 
we know how to generate sequences~\citep{feder1992universal} where $\Reg(\ftl)$ grows linearly with $n$, and so the $\sqrt{n}$ regret of $\cover$ becomes appealing. 
\bluet{This discussion indicates that $\ftl$'s poor worst-case performance occurs primarily on carefully constructed adversarial sequences, while achieving near-optimal performance on a broad class of naturally occurring stochastic processes.}

Now suppose instead that a fictitious `oracle' is told beforehand which of $\ftl$ or $\cover$ is better suited for the upcoming sequence $y^n$; the demand made by \emph{instance optimality} is that we try to be competitive against such an oracle \emph{on every sequence $y^n$}.
\begin{definition}[Instance Optimality]\label{def:instanceopt_binary}
A policy $\alg$ is instance optimal with respect to the regret of $\ftl$ and $\cover$ if there exists some universal $\gamma_n \geq 1$ such that for all $y^n \in \{0,1\}^n$:
%(independent of the sequence, and ideally, of $n$) 
\begin{align}
\Reg(\alg,y^n)
\leq \gamma_n \min\{\Reg(\ftl,y^n),\Reg(\cover,y^n)\}.
\label{eq:instanceopt_binary}
\end{align}
%for all $y^n \in \{0,1\}^n$. Equivalently, $\gamma_n$ is the competitive ratio achieved by $\alg$.
\end{definition}
We henceforth refer to $\gamma_n$ as the competitive ratio achieved by $\alg$; ideally we want this ratio to be a constant, i.e., $\gamma_n = O(1)$. 
This necessitates that on sequences where $\ftl$ gets a constant regret, then $\alg$ basically follows $\ftl$ throughout, while on sequences where $\ftl$ has high (in particular,~$\omega(\sqrt{n})$) regret, then $\alg$ follows $\cover$ in most rounds. 


The challenge in designing instance optimal algorithms is that the regret of any algorithm is a hindsight quantity that is not adapted to the natural filtration, i.e. it may not be possible to track $\Reg(\alg,y^n)$ for any $\alg$ from just the history $(y_1,y_2,\ldots,y_{t-1})$. 
One proxy is to use `corralling' policies that track each algorithm's loss instead, and try to be within a small additive loss compared to the better algorithm.
%leading to the idea of `corralling' policies~\cite{agarwal2017corralling,pacchiano2020model,Dann--Wei--Zimmert2023}, 
Treating both algorithms symmetrically results in the combined policy getting within  $O(\text{poly}(n))$ of the smaller of the two losses. This however can not ensure $\gamma_n=O(1)$:
for example, consider an i.i.d. sequence of Bernoulli$(0.1)$ bits, where $\ftl$ has lower regret than $\cover$. With high probability on any such sequence we have small $\Reg(\ftl,y^n)=O(1)$ and yet high loss $L_n(a^*,y^n)=\Theta(n)$; now any corralling algorithm (even a small loss one) must suffer $O(\text{poly}(n))$ regret, and hence $\omega(1)$ competitive ratio---\bluet{the final regret guarantee of 
\[
\min\{\Reg(\ftl, y^n), \Reg(\cover,y^n)\} + \Reg(\alg^{\text{corralling}}, y^n)
\]
is often dominated by the second term which can scale very differently from the first term}. This example also shows that achieving a constant factor guarantee with respect to the minimum of the two \emph{losses} does not translate to a constant factor guarantee with respect to the minimum of the two \emph{regrets}. We note though that treating the two algorithms asymmetrically does lead to stronger guarantees~\cite{sani2014exploiting}, though still not achieving $\gamma_n=O(1)$; we discuss this in more detail in~\cref{sec:relatedwork}.



% The only previous result which attains an instance-optimality guarantee is that of~\cite{de2014follow}, which proposes the $\textsf{AdaHedge}$ and $\textsf{FlipFlop}$ algorithms that carefully mix between $\ftl$ and a worst-case optimal algorithm called $\textsf{Hedge}$ in a data-driven way. 
% In our work, we propose a simple modular algorithm that takes as input any algorithm with a minimax guarantee and achieves $e/(e-1)$-instance optimal guarantees with respect to $\ftl$ and the worst case guarantee of the input algorithm via a simple reduction to the ski-rental problem. We also provide the first lower bound on the competitive ratio $\gamma$ for instance-optimality, which shows that our algorithm is close to optimal. \cycomment{There is some nuance here which makes our guarantee technically a little diff than the instance opt defn}

%\cycomment{Removing the footnote:\footnote{In particular ~\cite[Theorem $14$]{de2014follow} gives a guarantee of at most $5.64\cdot\Reg(\ftl)+4.64$ compared to $\ftl$, which both has a worse multiplicative constant, but also, is not instance-optimal due to the additional constant loss.} perhaps more suitable to put it in the discussion of the main result after it's stated in the body. Also because the competitive ratio constants are not really fair to compare since theirs is wrt the small loss bounds. Maybe here add a comment ours being black box?}

% More surprisingly, the guarantee is achieved via a simple approach, wherein we only switch at most once between $\ftl$ and $\wcalg$, and which only requires minimal ingredients.% -- monotone adapted regret estimates for $\ftl$ (which we provide in~\cref{lem:FTLRegret} for general online learning settings) and an upper bound on $\Reg(\wcalg)$.
% The allows us to extend the approach easily to a wide range of settings.


%\cycomment{The below paragraphs make more sense in the "model setup and def of instance optimality section" rather than "summary of results section". Probably need to fix the transition and reread for redundancies.}
We also consider a more general online learning setting where at the beginning of each round $t \in [n]$, a policy $\alg$ first plays an action $a_t \in \Ac$, following which, a loss function $\ell_t:\Ac\to [0,1]$ is revealed, resulting in a loss of $\ell_t(a_t)$. The \emph{regret} is defined according to:
\begin{align}
\label{eq:RegDefnGeneral}
    \Reg(\alg,\ell^n) = \textstyle\sum_{t=1}^n \ell_t(a_t) - \inf_{a \in \Ac} \textstyle\sum_{t=1}^n \ell_t(a).
\end{align} 
More generally, as in with binary prediction, we allow $\alg$ to play in round $t$ a measure $w_t\in\Delta_{\Ac}$ (i.e., play $\{w_t:\Ac\mapsto[0,1]|\sum_{a\in\Ac}w_t(a)=1\}$), resulting in an expected loss of $\sum_{a\in\Ac}w_t(a)\ell_t(a)$. For notational convenience, we henceforth use $(a_t,\ell_t(a_t))$ for the action/loss, and reserve use of expectations for randomness in the algorithm and/or sequence.

We want to understand when is it possible to attain instance optimality as in Eq.~\eqref{eq:instanceopt_binary} with respect to a given pair of algorithms. Ideally, we want the first to be optimal for stochastic instances, and the second to be minimax optimal; unfortunately however exact optimal policies are unknown except in simple settings. To this end, we make two amendments to our goal: First, for the stochastic optimal policy, we use $\ftl$; this is well defined in any online learning setting, and moreover, known to be optimal or near-optimal for a wide range of settings under minimal assumptions~\citep{kotlowski2018minimaxity}. Second, instead of the minimax policy, we use as reference \emph{any} policy $\wcalg$ which has a known worst case regret bound $g(n)$. With these modifications in place, we have the following objective.

%However, proving such a guarantee requires a tight characterization of the performance of both reference algorithms on every input sequence, which is not well understood for most algorithms. As a result, for general online learning, we consider the following weaker form of instance optimality
%with respect to the regret of $\ftl$ and a worst case regret bound $g(n)$.
\begin{definition} \label{def:instanceopt_wc}
Given \bluet{the follow-the-leader algorithm} $\ftl$ \bluet{which at time step $t$ selects action
\begin{align}\label{eq:FTLDefinition}
a_t = \arg \inf_{a \in \mathcal{A}} \sum_{i=1}^t \ell_i(a) 
\end{align}}
and any algorithm $\wcalg$ with 
\[
\sup_{\ell^n} \Reg(\wcalg, \ell^n) \le g(n),
\]
we say a policy $\alg$ is instance optimal with respect to the pair if there exists some universal $\gamma_n \geq 1$ such that for every sequence of losses $\ell^n$:
\begin{equation}\label{eq:instanceopt_wc}
 \Reg(\alg,\ell^n)\leq \gamma_n \min\{\Reg(\ftl,\ell^n),g(n)\}.
\end{equation}
%for all sequences of losses $\ell^n$. Equivalently, $\gamma$ is the competitive ratio achieved by $\alg$.
\end{definition}
%\cycomment{Is this possibly confusing in comparison to Definition 1? Does it make sense to define both formally? ...}
We note that Definition~\ref{def:instanceopt_wc} is an extension of the definition of instance-optimality for the binary setting (Definition~\ref{def:instanceopt_binary}) since in the binary setting, the worst-case regret $g(n) = \Reg(\cover,y_1^n) = \sqrt{\frac{n}{2\pi}}(1+o(1))$.
While strictly speaking the above guarantee is not truly instance-optimal in that we are comparing against a worst-case regret bound $g(n)$ for $\wcalg$ rather than its performance on the instance $\ell^n$, the two are the same if $\wcalg$ is minimax optimal and hence attaining equal regret on all loss sequences (as with the minmax algorithm $\cover$ for binary prediction). \bluet{It is interesting to define (and achieve) instance-optimality for a ``path" dependent worst-case bound; in Section~\ref{sec:SmallLoss}, we devise an algorithm that achieves a guarantee of the form $\Reg(\alg,\ell^n)\leq \gamma_n \min\{\Reg(\ftl,\ell^n),g(\ell^n)\}$, where $g(\ell^n)$ captures a certain measure of ``easiness" of the sequence $\ell^n$.}

\section{Summary of Results} \label{ssec:mainresults}
We propose the \emph{Switch via Monotone Adapted Regret Traces} ($\bob$) approach, which at a high level is a black-box way to convert design of instance-optimal policies into a simple \emph{optimal stopping} problem. 
The algorithm is extremely simple and formalizes the principle of {\em measured optimism}: ``assume the best'' by beginning with the opportunistic algorithm follow-the-leader $\ftl$, and ``prepare for the worst'' by switching to a given minimax optimal algorithm $\wcalg$ as needed when $\ftl$ collects too much loss.
Our analysis and guarantees build upon two key ingredients: first, owing to the additive structure of online learning problems, we have that the minimax guarantee $g(k)$ above holds over \emph{any} $k\in\mathbb{Z}$ and any (sub)sequence of $n$ loss functions; second, we show that $\ftl$ admits simple anytime regret estimator $\ftlsum_{\tau}$ (see Lemma~\ref{lem:FTLRegret}) which is \emph{monotone} and \emph{adapted} (i.e., a function only of historical data). Using these two observations, \bluet{we can analyze the task of minimizing regret using the competitive analysis framework from the `ski-rental' problem}~\cite{karlin1994competitive,borodin2005online}, as follows:
we play $\ftl$ up to some stopping time $\tau$, and then switch to $\wcalg$ for the remaining $n-\tau$ periods, resetting all losses to zero. This algorithm  incurs a \emph{total} regret bounded by $\ftlsum_{\tau} + g(n-\tau)$, and using ideas from competitive analysis, we get that there is a simple way to choose the stopping time $\tau$ to achieve an $e/(e-1) \approx 1.58$-competitive ratio guarantee with respect to the minimum between the regret of $\ftl$ and the worst case guarantee $g(n)$.

%abbreviated statement
{\bf Theorem \ref{thm:SkiRentalRegret} (simplified).}
{\em Let $\wcalg$ have worst-case regret $\sup_{\ell^n} \Reg(\wcalg, \ell^n) \le g(n)$
where $g(n)$ is some monotonic function of $n$. An instantiation of $\bob$ achieves  
    \begin{align}
\Reg(\bob, \ell^n)
&\le \frac{e}{e-1} \min\{\Reg(\ftl,\ell^n),g(n)\} + 1. \label{eq:2ApproxReg}
    \end{align}
}
A highlight of our approach is the surprising simplicity of the algorithm and analysis, despite the strength of the instance optimality guarantee. In particular, our approach is modular, allowing one to plug in any $\wcalg$ and corresponding worst case bound $g(n)$, thus letting us handle any online learning setting with known minimax bounds. This results in an entire family of instance optimal policies for settings such as predictions with experts and online convex optimization. Moreover, the approach is easy to extend to get more complex guarantees; as an example, if $\wcalg$ is designed to get low regret for benign (i.e., `small-loss') sequences $\ell^n$, then we show how to use $\bob$ as a subroutine and achieve an instance optimal guarantee with respect to the regret of $\ftl$ and a small loss regret bound. 

{\bf Corollary~\ref{cor:SqrtLStarInstantiation} (simplified).} 
\emph{Let $L^* \defeq \min_j \sum_{t=1}^n \ell_{tj}$. An instantiation of $\bob$ achieves}
    \begin{align*}
    %\label{eq:BOBWSqrtLstarRegDoubling}
        % \Reg(\pstar, \ell^n) &\le 2 \min\left\{\Reg(\ftl, \ell^n), 8 \sqrt{2L^* \log m} + \kappa\log\left(1+\frac{L^*}{\log m}\right)\log m + 8\sqrt{2}\log m \right\} \nonumber\\
        % &\qquad\qquad\qquad\qquad\qquad+ 2 \log\left(1+ \frac{L^*}{\log m}\right)
        \Reg(\bob, \ell^n) 
        &\le 2 \min\left\{ \Reg(\ftl,\ell^n), 10 \sqrt{2L^* \log m} \right\}
        + O(\log L^* \log m). 
    \end{align*}

Finally, studying instance optimality also lets us understand the fundamental limits of best-of-both worlds algorithms.
To this end, we provide a lower bound that shows our algorithm is nearly optimal in the competitive ratio. 
To the best of our knowledge, this is the first hardness result for best-of-both-worlds guarantees in online learning.

{\bf Theorem~\ref{thm:LowerBdCR}.} 
\emph{In the binary prediction setting, given \emph{any} online algorithm $\alg$, there exist sequences $y^n\in\{0,1\}^n$ such that:}
\begin{align*}
    \Reg(\alg,y^n)
    &\geq 1.43\min\left\{\Reg(\ftl,y^n),\Reg(\cover,y^n)\right\}.
\end{align*}
Recall again that in binary prediction, $\ftl$ achieves the optimal pseudoregret under i.i.d. inputs, while $\cover$ is the true minimax policy; thus, this is a fundamental lower bound on best-of-both-worlds guarantee in any online learning setting.



\bluet{The proof of Theorem~\ref{thm:LowerBdCR} uses Cover's characterization of achievable loss functions in binary prediction~\citep{Cover1966}. Any $\gamma$-competitive algorithm must satisfy a certain balance condition when applied to random sequences---we show this forces $\gamma \ge 1.4335$ by analyzing the expected performance on uniform random walks and computing the resulting integral constraints.}
